\capitulo{1}{Introducción}

El parkinson es una enfermedad crónica e irreversible que afecta a más de 150\,000 personas en toda España, siendo esta la segunda enfermedad neurodegenerativa más común después del Alzheimer~\cite{brainsci11081027}. Esta enfermedad se caracteriza por la perdida de neuronas dopaminérgicas lo que conlleva un desequilibrio en los circuitos neuronales que controlan el movimiento, lo cual puede provocar temblores, rigidez muscular, lentitud de movimiento (bradicinesia) y problemas de equilibrio y coordinación~\cite{Poewe_Seppi_Tanner_Halliday_Brundin_Volkmann_Schrag_Lang_2017}.
Pese a que no exista una cura, saber la evolución del paciente y su respuesta a los medicamentos y tratamientos puede ayudar enormemente tanto a pacientes como a médicos a adecuarse mejor a las circunstancias del paciente y con ello mejorar su calidad de vida.
 
La bradicinesia es uno de los síntomas más importantes a la hora de clasificar el grado de la enfermedad. Este síntoma se hace muy visible en un ejercicio conocido como test de golpeteo rápido de los dedos (o <<rapid finger tapping test>> en inglés). Este ejercicio consiste en separar y juntar repetidas veces los dedos índice y pulgar de forma rápida y constante.


\begin{figure}
    \centering
    \subfloat{\includegraphics[width=0.20\textwidth]{introduction/hand_open.jpeg}}
    \subfloat{\includegraphics[width=0.20\textwidth]{introduction/hand_closed.jpeg}}
    \subfloat{\includegraphics[width=0.20\textwidth]{introduction/hand_open.jpeg}}
    \subfloat{\includegraphics[width=0.20\textwidth]{introduction/hand_closed.jpeg}}
    \caption{Rapid finger tapping test}
    \label{fig:finger-tapping-test}
\end{figure}

En las últimas décadas, el campo de la inteligencia artificial ha experimentado un gran avance, con aplicaciones prometedoras en el ámbito de la salud. Específicamente, las técnicas de aprendizaje automático han demostrado su potencial para el diagnóstico y clasificación de enfermedades neurodegenerativas a partir de datos clínicos, genéticos y de neuroimagen~\cite{CardiacCare}.

En este proyecto se ha hecho uso de la inteligencia artificial para la creación de modelos capaces de clasificar el grado de bradicinesia en vídeos del <<rapid finger tapping test>>. Ver de manera individual la clasificación de cada vídeo puede no ser del todo útil para poder comprobar la evolución del paciente. Es por ello que se ha decidido crear, usando vídeos tomados de distintos días, una gráfica en la que se pueden ver las dos características de la bradicinesia: la amplitud y la lentitud. 

Este proyecto es la continuación de un Trabajo de Fin de Grado anterior llamado PADDEL~\cite{paddelRepo} desarrollado por Catalín Andrei Cacuci. La idea de este proyecto es pasar de un clasificador binario, que indicaba si el paciente presentaba o no la enfermedad de Parkinson, a un clasificador multiclase, el cual evalúa el grado de amplitud y lentitud del video. Además, como se indica anteriormente, se muestra una gráfica con la evolución del paciente.